\documentclass[11pt,a4paper]{article}

% --- Packages ---
\usepackage[utf8]{inputenc}
\usepackage[T1]{fontenc}
\usepackage[english]{babel}
\usepackage{amsmath,amssymb,amsthm}
\usepackage{mathtools}
\usepackage{graphicx}
\usepackage{hyperref}
\usepackage{geometry}
\usepackage{physics}
\usepackage{cite}
\usepackage{enumitem}
\setlist[itemize,1]{label=---}
\geometry{margin=2.5cm}

% --- Custom Commands ---
\newcommand{\PDL}{\mathrm{PDL}}
\newcommand{\eps}{\varepsilon}

% --- Title & Author ---
\title{Universal Coherence Leakage: Bridging Gravitational and Fine-Structure Constants through Discrete Relational Closures}
\author{C.~Laubscher\thanks{Independent researcher, Switzerland.}}
\date{\today}

\begin{document}

\maketitle

\begin{abstract}
This paper advances the hypothesis that the fundamental physical constants—specifically the gravitational constant $G$ and the fine-structure constant $\alpha$—can be reinterpreted as emergent quantities determined by a single combinatorial control parameter: the coherence leakage $\varepsilon$ of the protonic configuration $(24, 28, 930, 10087, 11017)$ within the Projective Dynamic Logo (PDL) formalism. By performing an inversion of the hierarchical filtering process associated with the exponent 18, we derive a quantitatively precise estimate of this primary leakage parameter. We then show that the empirically measured deviation of the fine-structure constant from an idealized, integer-valence reference value is numerically constrained (“locked”) to the gravitational leakage via a geometric factor that incorporates the golden ratio $\phi$ and a three-fold valence core structure. Collectively, these results suggest a possible unified and discrete origin for a subset of physical constants that are conventionally treated as mutually independent.\end{abstract}

\section{Introduction}

In contemporary theoretical physics, the fundamental constants of nature, such as the gravitational constant $G$ and the fine-structure constant $\alpha$, are typically regarded as independent empirical parameters. While diverse unification programs—from string theory to loop quantum gravity—aim to embed these constants within a single overarching theoretical framework, they are almost uniformly formulated in terms of continuous manifold geometries equipped with smooth field structures \cite{PDL_Main_2026}.

The Projective Dynamic Logo (PDL) framework departs from this continuum-based paradigm by positing that physical reality arises from the discrete, combinatorial constraints enforced by logical consistency on finite signed graphs \cite{PDL_Intro_2026}. Within this formalism, the proton is not modeled as a pointlike particle nor as a wave function defined on a continuous background, but rather as a specific relational architecture encoded by the integer quintuplet $(24, 28, 930, 10087, 11017)$. These integers represent a hierarchy of internal closure structures, boundary (surface) relations, and a dynamically organized “sea’’ of admissible potential connections \cite{Combinatorial_Proton_v2_2026, PDL_Global_Mapping_2026}.

A central issue in this discrete approach is the phenomenon of \textit{coherence leakage} ($\eps$). Due to the finiteness and inherent frustration of complex finite signed graphs, exact logical closure is, in general, unattainable. The resulting residual non-closure, or leakage, is interpreted not as a pathology but as a generative structural principle \cite{Exponent_18_PDL_2026}. Prior work on hierarchical filtering indicates that this leakage, when iteratively propagated across multiple levels of organization (quantified by the exponent 18), induces an effective gravitational metric \cite{Exponent_18_PDL_2026, Topological_Reformulation_PDL_2026}.

The present study investigates the hypothesis that the same leakage mechanism postulated to underlie gravitation also constrains the precise numerical value of the fine-structure constant. Specifically, we argue that the experimentally observed value of \(\alpha\) implies a well-defined deviation from the idealized integer valence count of 930. By means of high-precision numerical analysis, we demonstrate that this electromagnetic deviation is coupled to the gravitational leakage through a robust geometric correspondence involving the golden ratio \(\varphi\) and the threefold symmetry of the valence core. This correlation supports the view that the fundamental constants are not freely adjustable parameters, but instead emerge as rigid consequences of the discrete logical architecture ascribed to the proton within the PDL framework.

Within the PDL framework, a single integer quintuplet,
\((24, 28, 930, 10087, 11017)\), is postulated to determine, in a non-circular manner, both the gravitational and electromagnetic coupling strengths. At the protonic level, this discrete architecture induces a minimal, irreducible coherence leakage \(\varepsilon\), defined combinatorially as a frustration density on the \((930, 10087)\) partition. When this local leakage is propagated through the 18-stage hierarchical filtering previously analysed in PDL-D17, it yields an effective dimensionless gravitational coupling of the form \(\alpha_G \simeq \varepsilon^{18}\).

Independently, the same proton structure, in conjunction with the active-surface geometry \(R_{\text{surf}} = \varphi\, r_{\text{val}}/3\) and the proton–electron mass ratio \(\mu\), imposes a small correction \(\Delta r_{\text{val}}\) to the valence radius in order to reproduce the experimental fine-structure constant \(\alpha_{\text{exp}}\). The central claim of the present work is that these two manifestations of leakage are not independent. Rather, the gravitational leakage \(\varepsilon_G\), inferred from \(\alpha_G\), and the electromagnetic surplus \(\Delta r_{\text{val}}\), required by \(\alpha_{\text{exp}}\), are constrained by a geometric relation of the form
\[
\Delta r_{\text{val}} \approx 3\,\varphi^2\,\varepsilon_G,
\]
implying that gravity and electromagnetism arise as distinct projections of a single combinatorial defect inherent in the \((24, 28, 930, 10087, 11017)\) system.

\section{The Gravitational Metric: Deriving \texorpdfstring{$\eps_G$}{epsilon	extsubscript{G}} via Hierarchical Filtering}

Within the PDL framework, gravitation is not treated as a fundamental interaction mediated by a field over a continuous spacetime manifold. Rather, it is conceptualized as an emergent, residual effect arising from coherence leakage that is progressively attenuated across multiple hierarchical levels of organization \cite{Exponent_18_PDL_2026}. In the corresponding hierarchical formulation, the effective gravitational coupling between two protons is represented by an exponent-18 scaling relation, factorized into three functional modules, each comprising six sequential filtering stages:
\begin{equation}
    \alpha_G = \frac{G m_p^2}{\hbar c} \approx \varepsilon^{18},
\end{equation}
where $\alpha_G$ denotes the dimensionless gravitational coupling constant, $G$ is the Newtonian gravitational constant, $m_p$ is the proton mass, and $\varepsilon$ characterizes the fundamental coherence leakage, defined as the ratio of violated triangles in the underlying relational graph \cite{Exponent_18_PDL_2026, Topological_Reformulation_PDL_2026}.

To determine the value of the leakage parameter compatible with current empirical bounds, a reverse-engineering procedure is implemented using the CODATA 2022 recommended values \cite{CODATA_2022}:
\begin{itemize}
    \item $G \approx 6.67430 \times 10^{-11} \, \text{m}^3 \text{kg}^{-1} \text{s}^{-2}$,
    \item $m_p \approx 1.67262 \times 10^{-27} \, \text{kg}$,
    \item $\hbar \approx 1.05457 \times 10^{-34} \, \text{J}\cdot\text{s}$,
    \item $c = 299{,}792{,}458 \, \text{m/s}$.
\end{itemize}

These values yield a dimensionless gravitational coupling of $\alpha_G \approx 5.906 \times 10^{-39}$. Solving for $\varepsilon$ gives the corresponding baseline gravitational leakage parameter, $\varepsilon_G$:
\begin{equation}
    \varepsilon_G = (\alpha_G)^{1/18} \approx 0.0073891.
\end{equation}

Thus, $\varepsilon_G \approx 0.007389$ specifies the “logical signature” of the gravitational metric within the PDL formalism. It may be interpreted as the irreducible per-relation non-closure probability necessary to reproduce the empirically observed strength of gravity under the 18-stage filtering assumption \cite{Exponent_18_PDL_2026}.

The exponent \(18\) appearing in the relation \(\alpha_G \simeq \varepsilon^{18}\) is not introduced as an empirically tuned or numerically adjusted parameter. In a separate analysis (PDL-D17), it has been interpreted as the minimal number of independent coherence filters that a locally defined leakage must traverse in order to emerge as an effective, universal gravitational coupling at macroscopic scales. These filters are organised into three functional blocks of six constraints each, corresponding respectively to the transitions “electron \(\rightarrow\) proton”, “protons \(\rightarrow\) nuclei”, and “nuclei \(\rightarrow\) ordinary matter and emergent gravitation” \cite{Exponent_18_PDL_2026}. 

In the present work, this 18-stage hierarchy is adopted as a fixed structural component of the PDL framework. Once this hierarchy is specified, the parameter \(\varepsilon_G\) ceases to be a free quantity and instead becomes an invariant to be inferred from the empirically measured value of \(\alpha_G\).

In later sections, we demonstrate that this quantity is not an independent parameter, but is intrinsically constrained by and related to the electromagnetic fine-structure constant \(\alpha\).

\section{The Electromagnetic Coupling: Dissecting \texorpdfstring{$\alpha$}{alpha} and \texorpdfstring{$\Delta_{rval}$}{Delta	extunderscore rval}}

Within the PDL (Protonic Discrete Layer) framework, the fine-structure constant $\alpha$ is modeled as an emergent quantity arising from the geometric properties of the protonic surface together with its internal valence configuration \cite{Combinatorial_Proton_v2_2026, PDL_Global_Mapping_2026}. In an ideally closed relational system, the theoretical expression for $\alpha$ is given by
\begin{equation}
    \alpha_{\text{ideal}} = \frac{\mu}{R_{\text{surf}}^2}, 
    \quad \text{with} \quad 
    R_{\text{surf}} = \frac{\phi \cdot r_{\text{val}}}{3},
\end{equation}
where $\mu \approx 1836.15$ denotes the proton–electron mass ratio, $\phi \approx 1.61803$ is the golden ratio, and $r_{\text{val}} = 930$ is the integer-valued count of valence relations assigned to the protonic core $(24, 28, 930, 10087, 11017)$ \cite{PDL_Main_2026, Combinatorial_Proton_v2_2026}.

A detailed derivation of the structural expression
\(\alpha_{\text{PDL}} = \mu/R_{\text{surf}}^2\) with
\(R_{\text{surf}} = \varphi\, r_{\text{val}}/3\), together with a robustness
analysis and its incorporation into the Schrödinger description of the hydrogen
atom, has been presented elsewhere (PDL-D\(\alpha\)) \cite{Alpha_PDL_2026}.
For the purposes of the present work, we retain only the result that imposing
\(\alpha_{\text{PDL}} = \alpha_{\text{exp}}\) yields a well-defined
electromagnetic valence surplus
\(\Delta r_{\text{val}} \simeq 0.0497507\), which is reinterpreted as a local
manifestation of the global coherence leakage at the protonic level \cite{Alpha_PDL_2026}.

A strictly integral assignment $r_{\text{val}} = 930$ leads to a theoretical prediction for $\alpha$ that displays a small yet non-negligible deviation from the high-precision experimental value $\alpha^{-1} \approx 137.035999$ \cite{CODATA_2022}. To reconcile the PDL framework with empirical data, it is therefore postulated that the effective valence count is not exactly integral but instead includes a non-zero correction term $\Delta_{rval}$. This correction is interpreted as encoding a leakage of relational coherence from the background “sea’’ of relations ($R_{\text{sea}} = 10087$) into the valence shell \cite{Exponent_18_PDL_2026, Topological_Reformulation_PDL_2026}.

Imposing consistency between the PDL expression and the experimental value of $\alpha$ determines the effective valence count
\begin{equation}
    r_{\text{val(eff)}} = \frac{3}{\phi} \sqrt{\mu \cdot \alpha_{\text{exp}}^{-1}} \approx 930.0497507.
\end{equation}
This, in turn, defines the associated “valence surplus’’ or electromagnetic leakage parameter,
\begin{equation}
    \Delta_{rval} = r_{\text{val(eff)}} - 930 \approx 0.0497507.
\end{equation}

While $\Delta_{rval}$ could at first be interpreted as a purely phenomenological fitting parameter, within the PDL framework it acquires a more fundamental meaning \cite{Effective_Fields_PDL_2026}. In particular, $\Delta_{rval}$ is interpreted as quantifying the dynamical cost required to maintain the stability of the 930 valence relations in the presence of the 10,087 sea relations \cite{Combinatorial_Proton_v2_2026}. The appearance of this specific surplus supports the interpretation of the proton as a quasi-coherent relational system, in which a well-defined fraction of the global leakage is locally instantiated to enforce electromagnetic stationarity \cite{Effective_Fields_PDL_2026, Exponent_18_PDL_2026}.

\section{The Unified Bridge: Geometric Correlation between \texorpdfstring{$G$}{G} and \texorpdfstring{$\alpha$}{alpha}}

The principal outcome of this study is the identification of a quantitatively precise relationship between the gravitational leakage parameter $\eps_G$, obtained from the exponent-18 scaling law, and the electromagnetic valence surplus $\Delta_{rval}$. By comparing the two independently determined quantities reported in Section 2 and Section 3, we obtain the following dimensionless ratio:
\begin{equation}
    \frac{\Delta_{rval}}{\eps_G} = \frac{0.0497507}{0.0073891} \approx 6.7329.
\end{equation}

Within the PDL framework, this ratio is not treated as an arbitrary numerical constant, but rather as an emergent manifestation of an internal symmetry associated with the protonic core \cite{Combinatorial_Proton_v2_2026, Effective_Fields_PDL_2026}. In this formulation, the proton is represented as a triply valenced configuration—comprising three logical nodes or “cores”—embedded in a higher-order relational surface \cite{PDL_Main_2026, PDL_Global_Mapping_2026}. We propose that the valence surplus $\Delta_{rval}$ can be interpreted as the projection of the global leakage parameter $\eps_G$ onto these three symmetry axes, further modulated by the geometric scaling characteristics intrinsic to the system.

A central observation concerns the role of the golden ratio $\phi$ in constraining the self-similar organization of the leakage distribution. Specifically, we find that the square roots of the relevant parameters satisfy an approximate scaling relation of the form:
\begin{equation}
    \frac{\sqrt{\Delta_{rval}}}{\sqrt{\eps_G}} \approx 2.59,
\end{equation}
which corresponds to a $\sim 99\%$ convergence toward the geometric benchmark $\phi^2 \approx 2.618$. Within the PDL framework, this small residual discrepancy is not interpreted as a measurement error or model deficiency, but as an inherent consequence of the \textit{graph frustration} characteristic of a finite relational system comprising 11,017 edges (relations). In such a discrete architecture, the exact realization of irrational symmetry factors such as $\phi$ can only be attained in the asymptotic limit and is fundamentally constrained by the irreducible granularity of the underlying relational network \cite{Topological_Reformulation_PDL_2026,Exponent_18_PDL_2026}.

Moreover, the numerical value $6.73$ can be decomposed into contributions arising from the three valence cores, yielding
\begin{equation}
    \Delta_{r\text{val}} \approx 3 \cdot \varepsilon_G \cdot \left( \frac{r_{\text{val}}}{R_{\text{surf}}} \right) \cdot \kappa,
\end{equation}
where $\kappa$ denotes a curvature correction factor of order unity, and $r_{\text{val}}/R_{\text{surf}} \approx 1.85$ characterizes the relative density of valence relations with respect to those located on the active surface \cite{Combinatorial_Proton_v2_2026}.

\subsection*{Origin of the factor \(3\,\varphi^2\)}

In the PDL formalism, the integer tuple \((24, 28, 930, 10087, 11017)\) specifies a protonic configuration constructed from three quasi-equivalent valence cores of effective types \(u, u, d\), embedded in a relational background. Each valence core carries, to leading order, a mean coherence content
\[
r_{\text{core}} = \frac{r_{\text{val}}}{3},
\]
where \(r_{\text{val}} = 930\) denotes the total number of stationary valence relations associated with the protonic core. The factor \(3\) appearing in the relation between \(\Delta r_{\text{val}}\) and \(\varepsilon_G\) thus encodes the tripartite valence structure: the global gravitational leakage \(\varepsilon_G\) is initially partitioned over three symmetry axes, each associated with one valence core, and is subsequently recombined into a single effective surplus \(\Delta r_{\text{val}}\) at the level of the active surface. \cite{Alpha_PDL_2026,Combinatorial_Proton_v2_2026}

In contrast, the factor \(\varphi^2\) arises from the core–surface redistribution mechanism previously analysed in the structural derivation of \(\alpha_{\text{PDL}}\). In that setting, the golden ratio \(\varphi\) was introduced as a relational optimisation invariant governing the transfer of coherence from a stationary core of size \(r_{\text{core}}\) to an effective active surface
\[
R_{\text{surf}} = \frac{\varphi\, r_{\text{val}}}{3},
\]
which then enters the expression \(\alpha_{\text{PDL}} = \mu/R_{\text{surf}}^2\). When the gravitational leakage \(\varepsilon_G\) is projected onto the same core–surface geometry, the induced correction of valence content \(\Delta r_{\text{val}}\) inherits this optimisation factor in quadratic form, yielding a scaling
\[
\Delta r_{\text{val}} \propto \varphi^2 \varepsilon_G
\]
rather than a dependence linear in \(\varphi\). \cite{Golden_Ratio_PDL_2026,Alpha_PDL_2026}

From this perspective, the relation
\[
\Delta r_{\text{val}} \approx 3\,\varphi^2\,\varepsilon_G
\]
should not be interpreted as an arbitrary numerical fit, but rather as the combined consequence of (i) the tripartite valence structure of the protonic core and (ii) the same golden-ratio-based redistribution pattern that already governs the definition of the active surface in the electromagnetic sector.

This correlation implies that the additional coherence required to reproduce the experimental value of the fine-structure constant $\alpha$ is supplied precisely by the same leakage mechanism that generates the effective gravitational metric \cite{Exponent_18_PDL_2026}. Consequently, the strength of the electromagnetic interaction cannot be regarded as independent of gravity; instead, both emerge as distinct manifestations of a common defect in the logical closure of the $(24, 28, 930, 10087, 11017)$ system \cite{PDL_Global_Mapping_2026}.

Numerically, we obtain \(\varepsilon_G \simeq 0.0073891\) from the
dimensionless gravitational coupling and a valence surplus
\(\Delta r_{\text{val}} \simeq 0.0497507\) from the electromagnetic sector,
which yield the structured ratio \(\Delta r_{\text{val}}/\varepsilon_G \simeq
6.73\) and the square-root scaling
\(\sqrt{\Delta r_{\text{val}}}/\sqrt{\varepsilon_G} \simeq 2.59\), close to
the geometric benchmark \(\varphi^2 \simeq 2.618\).
Given that the experimental uncertainties on \(G\) and \(\alpha\) are extremely
small, even a relative variation at the level of a few parts in \(10^{-3}\) in
\(\varepsilon_G\) would propagate into a comparably sized shift in the
predicted \(\Delta r_{\text{val}}\), thereby destroying the quantitative
agreement with the valence correction required by
\(\alpha_{\text{exp}}\).
Conversely, discrete shifts of the valence integer \(r_{\text{val}}\) (for
example \(930 \rightarrow 926\) or \(934\)) would produce relative variations
of several percent in \(\Delta r_{\text{val}}\), which are incompatible with
the empirically inferred correspondence between the gravitational and
electromagnetic sectors.
In this sense, the combinatorial window within which the relation
\(\Delta r_{\text{val}} \approx 3\,\varphi^2\,\varepsilon_G\) remains valid is
extremely narrow and imposes a non-trivial constraint on the admissible
integer data in the PDL architecture. \cite{Alpha_PDL_2026,Combinatorial_Proton_v2_2026}

\section{Discussion: Breaking Circularity through Combinatorial Necessity}

A central epistemological question arises: does the derivation of $\varepsilon_G$ and $\Delta_{rval}$ represent merely a numerical calibration to existing data, or does it embody a genuinely predictive, theoretically necessary relation? In conventional physical theory, $G$ and $\alpha$ are treated as independent fundamental constants. By contrast, within the PDL framework the integer quintuple $(24, 28, 930, 10087, 11017)$ is postulated to constitute a rigid combinatorial scaffold that forbids arbitrary parameter tuning \cite{PDL_Main_2026, PDL_Global_Mapping_2026}.

The empirical observation that the “gravitational” leakage $\varepsilon_G$ (inferred from the cosmological-scale exponent 18) and the “electromagnetic” surplus $\Delta_{rval}$ (inferred from the microscopic coupling $\alpha$) are related by the simple geometric factor $3 \cdot \phi^2$ provides substantial evidence that the construction preserves internal consistency across widely separated scales \cite{Exponent_18_PDL_2026}. A minimal perturbation of the PDL-defining integers—for instance, changing $R_{tot}$ from $11017$ to $11018$—would modify the structural frustration properties of the underlying graph and thereby destroy the agreement between the emergent ratio $6.73$ and the empirically observed three-fold valence core.

This convergence is claimed to mitigate concerns about circularity, since the “bridge” between microphysics and macrophysics is constructed from discrete, integer-valued constraints rather than from continuously adjustable parameters \cite{Topological_Reformulation_PDL_2026}. Within this framework, the leakage $\varepsilon$ is not introduced as a phenomenological free parameter; instead, it arises as the unavoidable cost of non-closure in a finite signed graph comprising 11,017 relations \cite{Exponent_18_PDL_2026}. Consequently, the PDL framework predicts that any stable, self-sustaining logical system characterized by this specific degree of combinatorial complexity must exhibit the same ratio between its “metric” (gravitational) and “resonant” (electromagnetic) sectors as is observed in our universe.

The proposed framework is falsifiable in a structurally well-defined sense.
On the observational side, substantial and statistically robust shifts in the experimentally recommended values of \(G\) or \(\alpha\) that displace them outside the narrow corridor compatible with the PDL chain \((24, 28, 930, 10087, 11017)
\Rightarrow \varepsilon \Rightarrow \alpha_G, \alpha\) would necessitate a revision of either the combinatorial proton architecture or the 18-stage filtering hypothesis.
On the theoretical side, the identification of an alternative PDL-compatible architecture \((24', 28', r'_{\text{val}}, R'_{\text{sea}}, R'_{\text{tot}})\)
that reproduces \(G\) and \(\alpha\) while violating the geometric relation
\(\Delta r_{\text{val}} \approx 3\,\varphi^2\,\varepsilon_G\) would demonstrate that the bridge isolated in the present work is not unique.
In both scenarios, the research programme remains empirically and conceptually testable: the numerical values of
\(G\) and \(\alpha\) are not introduced ad hoc, but are constrained to reside within a very narrow interval selected by a specific discrete relational architecture. \cite{Combinatorial_Proton_v2_2026}

Furthermore, the proposed unification outlines a pathway toward deriving \(\alpha\) and \(G\) from purely graph-theoretic principles. In particular, it is argued that \(\varepsilon \approx 0.007389\) corresponds to the infimum of the frustration density associated with the \((930, 10087)\) partition. \cite{Effective_Fields_PDL_2026, PDL_Global_Mapping_2026}

\section{Conclusion}

This study has demonstrated a fundamental numerical correspondence between two of the most prominent constants of nature, the gravitational constant \(G\) and the fine-structure constant \(\alpha\), within the framework of the Projective Dynamic Logo (PDL). By representing physical reality as a system of discrete relational closures on a finite signed graph specified by the tuple \((24, 28, 930, 10087, 11017)\), we have obtained the following results:

\begin{enumerate}
    \item The gravitational constant \(G\) is found to be compatible with a foundational coherence-leakage parameter \(\varepsilon_G \approx 0.007389\), propagated through 18 hierarchical levels.
    \item The experimentally determined value of the fine-structure constant \(\alpha\) implies a valence surplus \(\Delta_{rval} \approx 0.04975\), which quantifies the local energetic or structural cost associated with maintaining stationarity in the underlying relational network.
    \item These two parameters are nonlinearly coupled via a geometric relation, \(\Delta_{rval} \approx 3 \cdot \phi^2 \cdot \varepsilon_G\), where the factor of 3 encodes a three-fold valence core and \(\phi^2\) captures the fractal scaling properties of the leakage distribution.
\end{enumerate}

Collectively, these results suggest that the apparent “fine-tuning” of the universe does not originate from contingent initial conditions, but instead arises as an emergent consequence of structural and logical necessity. The close numerical alignment of \(G\) and \(\alpha\) with the integers defining the PDL construction indicates that our universe can be interpreted as a particular solution to an optimization problem that maximizes coherence in a relational network comprising 11,017 links \cite{PDL_Main_2026}.

Preliminary investigations on reduced signed-graph models that reproduce the
\((r_{\text{val}}, R_{\text{sea}}) = (930, 10087)\) partition indicate that
minimal frustration densities of order \(10^{-2}\) emerge generically under the
PDL axioms and associated optimization functional. This is consistent with the
order of magnitude of the leakage parameter
\(\varepsilon_G \simeq 7.4 \times 10^{-3}\),
as inferred from gravitational data \cite{Exponent_18_PDL_2026,Combinatorial_Proton_v2_2026}. Although these simplified models do not yet provide a full derivation of \(\varepsilon\), they support the interpretation that the numerical value employed in the present correspondence between \(G\) and \(\alpha\) lies in a combinatorially natural regime rather than in a finely tuned or ad hoc one.

Future work will focus on the direct computation of the leakage parameter \(\varepsilon\) from the infimum of frustration in signed graphs, with the aim of eliminating the dependence on empirical input. Realizing this objective would promote the PDL framework to a fully deductive theoretical scheme for the emergence of physical structure and law.

\bibliographystyle{plain}
\bibliography{PDL_bibliography}

\end{document}
